% !TEX root = ../../main.tex
% ============================================================================
% Chapter: Universality of the Sugawara denominator (k + h^v)
%
% Russian-school lossless adversarial attack + Platonic reconstitution.
% Drinfeld + Feigin-Frenkel + Faddeev-Takhtajan style.
%
% This chapter STRENGTHENS the convention paragraphs already present at
% standalone/chiral_chern_weil.tex:474
% standalone/holographic_datum.tex:636
% standalone/three_parameter_hbar.tex
% standalone/virasoro_r_matrix.tex
% from a pointwise bridge identity to a UNIVERSALITY theorem. The
% pre-existing paragraphs correctly dismiss the naive rational identity
% k Omega_tr = Omega_KZ / (k + h^v);
% the structural content behind that
% dismissal: the denominator (k + h^v) is a universal invariant of the
% affine Kac-Moody chiral datum, appearing in the Sugawara conformal
% weight, the Sugawara central charge, the trace-form shadow invariant
% kappa(V_k(g)), the KZ connection, and the Feigin-Frenkel critical-level
% locus; the trace-form and KZ r-matrix presentations are two gauge
% representatives of the SAME operadic datum, exchanged by a
% level-dependent change of basis of the Casimir.
%
% Guiding principle: strengthening only. Every claim is scoped to the
% convention (trace-form / KZ) in which it is stated. No bare kappa;
% label prefixes match environment; r-matrix writes carry level prefix
% and k=0 check.
% ============================================================================

\providecommand{\fg}{\mathfrak{g}}
\providecommand{\fh}{\mathfrak{h}}
\providecommand{\fz}{\mathfrak{z}}
\providecommand{\cA}{\mathcal{A}}
\providecommand{\cC}{\mathcal{C}}
\providecommand{\cO}{\mathcal{O}}
\providecommand{\Omegatr}{\Omega_{\mathrm{tr}}}
\providecommand{\OmegaKil}{\Omega_{\mathrm{Kil}}}
\providecommand{\OmegaKZ}{\Omega_{\mathrm{KZ}}}
\providecommand{\rtr}{r_{\mathrm{tr}}}
\providecommand{\rKZ}{r_{\mathrm{KZ}}}
\providecommand{\hv}{h^{\vee}}
\providecommand{\kch}{\kappa_{\mathrm{ch}}}

\chapter{Universality of the Sugawara denominator $k + \hv$}
\label{chap:chern-weil-level-shift-platonic}
\index{Sugawara!denominator $k + h^{\vee}$|textbf}
\index{level shift!universality|textbf}
\index{trace-form convention!vs.\ KZ convention|textbf}
\index{Feigin--Frenkel!critical level|textbf}

\ChapterIntro{The affine Kac--Moody chiral algebra $V_k(\fg)$ carries
two natural classical $r$-matrix presentations. In the \emph{trace-form}
convention,
\[
 \rtr(z) \;=\; \frac{k\,\Omegatr}{z},
\]
where $\Omegatr = \sum_a t^a \otimes t^a$ is the Casimir built from an
orthonormal basis for the normalised invariant form (long roots squared
length~$2$). In the \emph{Kazhdan--Lusztig} (KZ) convention,
\[
 \rKZ(z) \;=\; \frac{\OmegaKZ}{(k + \hv)\,z},
\]
where $\OmegaKZ=\Omegatr$ is the same normalised inverse-form Casimir
written in conformal-block notation. If $\OmegaKil$ denotes the inverse
tensor of the Killing form, then
$\OmegaKil=(2\hv)^{-1}\Omegatr$ and
$\OmegaKZ=2\hv\,\OmegaKil$. These two formulas are \emph{not} equal
as rational functions of~$k$: at $k = 0$ the left side vanishes while
the right side is finite and non-zero. This chapter makes the
structural content behind that distinction precise. The denominator
$k + \hv$ is a universal
invariant of the chiral datum, appearing simultaneously in the
Sugawara energy--momentum tensor, the Sugawara central charge, the
trace-form shadow invariant, the KZ connection, and the Feigin--Frenkel
critical-level locus. The trace-form and KZ presentations are not
generic tensor-equal representatives; they are the central-extension
and Sugawara-Ward projections of the same affine operadic datum. The
convention discipline is universal: every $r$-matrix write declares
its convention and exhibits the appropriate $k = 0$ check.}

\section{Setup: affine Kac--Moody at level $k$}
\label{sec:chern-weil-setup}

Fix a simple Lie algebra $\fg$ over $\mathbb{C}$ with dual Coxeter
number $\hv$ and rank $r$. Let $\kappa_{\mathrm{Kil}}$ be the Killing
form and let $(\ ,\ )$ denote the normalised invariant form determined
by the convention $(\theta, \theta) = 2$ for the highest root
$\theta$. The two invariant forms are related by
\begin{equation}\label{eq:killing-trace-relation}
 \kappa_{\mathrm{Kil}}(x, y) \;=\; 2 \hv\, (x, y),
 \qquad x, y \in \fg.
\end{equation}

Let $\{t^a\}_{a=1}^{\dim\fg}$ be an orthonormal basis for $(\ ,\ )$
and let $\{T^a\}$ be an orthonormal basis for
$\kappa_{\mathrm{Kil}}$. The two Casimirs are
\begin{equation}\label{eq:casimir-two-forms}
 \Omegatr \;=\; \sum_a t^a \otimes t^a,
 \qquad
 \OmegaKil \;=\; \sum_a T^a \otimes T^a,
 \qquad
 \OmegaKil \;=\; \frac{1}{2 \hv}\,\Omegatr.
\end{equation}
Explicitly, $T^a = (2\hv)^{-1/2} t^a$ gives an orthonormal basis
for $\kappa_{\mathrm{Kil}}$; squaring the normalisation factor recovers
the prefactor $(2\hv)^{-1}$ in~\eqref{eq:casimir-two-forms}. The KZ
Casimir used below is
\begin{equation}\label{eq:kz-casimir-normalization}
 \OmegaKZ := \Omegatr = 2\hv\,\OmegaKil.
\end{equation}
Thus the KZ numerator is not the Killing inverse tensor itself; it is
the normalised inverse-form tensor expressed in the notation standard
for conformal-block connections.

The affine Kac--Moody algebra at level $k$ is the vertex algebra
$V_k(\fg)$ with strong generators $J^a(z)$ and singular OPEs
\begin{equation}\label{eq:affine-km-ope}
 J^a(z)\,J^b(w) \;=\;
 \frac{k\,\delta^{ab}}{(z - w)^2}
 \;+\;
 \frac{f^{ab}{}_c\,J^c(w)}{z - w}
 \;+\;\mathrm{regular},
\end{equation}
where $f^{ab}{}_c$ are the structure constants in the $\{t^a\}$
basis. The level $k$ multiplies the Killing-normalised central
extension after the orthonormal-basis translation: under $t^a \mapsto
T^a = (2\hv)^{-1/2} t^a$, the OPE~\eqref{eq:affine-km-ope} becomes
\[
 J^a_{\mathrm{KZ}}(z)\,J^b_{\mathrm{KZ}}(w)
 \;=\;
 \frac{(k / 2\hv)\,\delta^{ab}_{\mathrm{Kil}}}{(z - w)^2}
 \;+\;\cdots,
\]
so in the KZ convention the central extension appears with effective
level $k/(2\hv)$.

\begin{definition}[r-matrix convention labels]
\label{def:rmatrix-convention}
% r-matrix check:
% family: affine KM; both conventions.
% convention (trace-form): r(z) = k * Omega_tr / z; level k MANDATORY.
% r(z)|_{k=0} = 0.
% convention (KZ): r(z) = Omega_KZ / ((k+h^v) z); level denominator MANDATORY.
% r(z)|_{k=0} = Omega_KZ / (h^v z), non-zero for non-abelian g.
% KZ r(z) does not vanish at k=0 because the Lie bracket persists;
% this is the physically correct statement. Critical level:
% r(z)|_{k=-h^v} diverges.
% Source: landscape_census.tex for kappa; this chapter for the bridge.
We fix once and for all the following two classical $r$-matrix
presentations of the affine Kac--Moody chiral algebra:
\begin{enumerate}
 \item \emph{Trace-form convention.} $\rtr(z) = k\,\Omegatr / z$,
 where $\Omegatr$ is built from an orthonormal basis for the
 normalised invariant form $(\ ,\ )$. The level $k$ appears as an
 explicit numerator prefix. At $k = 0$: $\rtr = 0$.
 \item \emph{Kazhdan--Lusztig} (KZ) \emph{convention.}
 $\rKZ(z) = \OmegaKZ / ((k + \hv)\,z)$, where $\OmegaKZ$ is built
 from the normalised form and equals $2\hv$ times the Killing
 inverse tensor. The level appears in the denominator as
 $k + \hv$. At $k = 0$:
 $\rKZ = \OmegaKZ / (\hv\,z)$, which is non-zero for non-abelian
 $\fg$ (the Lie bracket persists through the Casimir).
\end{enumerate}
\end{definition}

\begin{remark}[bare-$r$ regime]
\label{rem:ap126-universal}
An $r$-matrix presentation that writes $\Omega / z$ with neither a $k$
numerator prefix nor a $(k + \hv)$ denominator is under-specified: it
chooses neither convention and so cannot be checked against either
$k = 0$ boundary. The regex
\[
 \mathtt{r(z)\backslash s*=\backslash s*\backslash\backslash Omega\backslash s*/\backslash s*z}
\]
catches this violation; every such bare write must be augmented into
$\rtr$ or $\rKZ$ form before the file is accepted. The universality
theorem below explains why both augmentations are legitimate only after
their extraction channel has been declared.
\end{remark}

\section{Universality of the level shift}
\label{sec:universality}

The assignment $k \mapsto k + \hv$ is not a change of basis that turns
$k\,\Omegatr$ into $\OmegaKZ/(k+\hv)$. It is the unique denominator
forced by the Sugawara normalisation, and hence by the KZ equation,
while $k\,\Omegatr$ is the central-extension part of the affine OPE.

\begin{theorem}[Level-shift universality with convention separation]
\label{thm:level-shift-universality}
\ClaimStatusProvedHere
Let $\fg$ be simple and $k \neq -\hv$. Consider the space of classical
$r$-matrix numerators and KZ connection coefficients associated with
$V_k(\fg)$ in the normalised invariant-form convention. Then:
\begin{enumerate}[label=\textup{(\roman*)}]
\item the trace-form collision residue is
\[
 \rtr(z)=\frac{k\,\Omegatr}{z};
\]
\item the KZ conformal-block connection has simple-pole coefficient
\[
 \rKZ(z)=\frac{\OmegaKZ}{(k+\hv)z},
 \qquad \OmegaKZ=\Omegatr=2\hv\,\OmegaKil;
\]
\item no $k$-independent invariant-form rescaling identifies these two
 tensors for generic~$k$;
\item the denominator $k+\hv$ is uniquely forced by the requirement that
 the Sugawara field make each current $J^a$ a primary field of
 conformal weight~$1$ and by the critical-level condition that the
 Sugawara construction degenerate precisely at $k=-\hv$.
\end{enumerate}
\end{theorem}

\begin{proof}
The affine OPE in the basis orthonormal for the normalised form is
\[
 J^a(z)J^b(w)
 =
 \frac{k\,\delta^{ab}}{(z-w)^2}
 +\frac{f^{ab}{}_cJ^c(w)}{z-w}
 +\cdots .
\]
The double-pole term is exactly the trace-form collision residue
$k\,\Omegatr/z$, proving~(i).

The invariant-form conversion is independent of~$k$:
equations~\eqref{eq:casimir-two-forms}--\eqref{eq:kz-casimir-normalization}
give $\OmegaKil=(2\hv)^{-1}\Omegatr$ and
$\OmegaKZ=\Omegatr=2\hv\,\OmegaKil$. This proves the tensor
normalisation in~(ii).

It remains to derive the denominator. Write a possible Sugawara field
as
\[
 T_a(z)=a\sum_b :J^bJ^b:(z).
\]
Contracting $T_a(z)$ with $J^c(w)$ using the affine OPE gives the
singular part
\[
 T_a(z)J^c(w)
 =
 \frac{2a(k+\hv)J^c(w)}{(z-w)^2}
 +\frac{2a(k+\hv)\partial J^c(w)}{z-w}
 +\cdots .
\]
The $k$ term comes from the central double-pole contraction, and the
$\hv$ term is the adjoint quadratic-Casimir contribution
$\sum_{a,b} f^{ca}{}_b f^{ab}{}_d = 2\hv\,\delta^c_d$ in the present
normalisation, with the factor~$2$ already absorbed by the two Wick
contractions. The condition that $J^c$ be primary of weight~$1$ forces
\[
 2a(k+\hv)=1,
 \qquad\text{so}\qquad
 a=\frac{1}{2(k+\hv)}.
\]
The same coefficient is the Ward-identity coefficient in the KZ
connection, proving~(ii) and the uniqueness part of~(iv).

Finally, if one tries to identify $\rtr$ and $\rKZ$ by a
$k$-independent invariant-form rescaling $\OmegaKZ=\lambda\Omegatr$,
then equality would require
\[
 k=\frac{\lambda}{k+\hv}
 \qquad\text{for all }k,
\]
or $k(k+\hv)=\lambda$ for all~$k$, impossible for constant
$\lambda$. Thus (iii) holds. The denominator vanishes exactly at
$k=-\hv$, the Feigin--Frenkel critical level, so the critical-locus
condition also forces the same shift.
\end{proof}

\begin{remark}[Change of basis vs.\ rational identity]
\label{rem:basis-vs-identity}
The equation
\[
 k \cdot \Omegatr \;\neq\; \OmegaKZ / (k + \hv)
\]
FAILS as a rational identity in $k$: at $k = 0$ the left side vanishes,
the right side is $\OmegaKZ / \hv \neq 0$. The correct reading is that
the two sides are different projections of the same affine OPE. The
left side records the central-extension collision residue. The right
side records the Ward-identity coefficient of the KZ connection after
the Sugawara field has been normalised by $1/(2(k+\hv))$. The
invariant-form conversion is only
	$\OmegaKZ=\Omegatr=2\hv\,\OmegaKil$; it does not supply any
	$k$-dependent tensor equality. The level-shift theorem records the
	positive statement: the trace-form central residue and the KZ
	Ward-identity coefficient are different projections of the same affine
	OPE, separated by the Sugawara denominator.
\end{remark}

\section{Universal occurrence of $k + \hv$}
\label{sec:universal-occurrence}

The level shift $k \mapsto k + \hv$ is not an accident of the
$r$-matrix bridge; the denominator $k + \hv$ appears universally in
every Sugawara-related quantity attached to $V_k(\fg)$.

\begin{proposition}[Universal occurrence of $k + \hv$]
\label{prop:shift-appears-universally}
\ClaimStatusProvedHere
Let $\fg$ be simple and $k \neq -\hv$. The quantity $k + \hv$ appears
as the Sugawara denominator in each of the following invariants of
$V_k(\fg)$:
\begin{enumerate}
 \item \emph{Sugawara energy-momentum tensor.}
 \[
 T(z) \;=\; \frac{1}{2(k + \hv)}
 \sum_{a} :\!J^a(z)\,J^a(z)\!:.
 \]
 \item \emph{Sugawara central charge.}
 \[
 c(V_k(\fg)) \;=\; \frac{k\,\dim \fg}{k + \hv}.
 \]
 \item \emph{Trace-form chiral shadow invariant.}
 \[
 \kch(V_k(\fg)) \;=\; \frac{\dim(\fg)\,(k + \hv)}{2\,\hv}
 \qquad\text{[trace convention]},
 \]
 equivalently, in KZ convention,
 $\kch(V_k(\fg))/(k + \hv) = \dim(\fg)/(2\hv)$, a
 level-independent rational number.
 \item \emph{KZ connection.} The genus-$0$ shadow connection on the
 configuration space $\mathrm{Conf}_n(\mathbb{C})$ is
 \[
 \nabla^{\mathrm{hol}}_{0,n}
 \;=\; d
 \;-\; \frac{1}{k + \hv}\sum_{i < j}
 \frac{\Omega_{\mathrm{KZ},\,ij}\,dz_{ij}}{z_{ij}}.
 \]
 \item \emph{Feigin-Frenkel critical level.} The locus where the
 Sugawara construction degenerates is $k = -\hv$, equivalently the
 zero locus of the universal denominator $k + \hv$.
\end{enumerate}
The Sugawara and KZ quantities are ill-defined precisely at
$k = -\hv$; the trace-form shadow invariant is instead regular there
and vanishes. In every displayed formula the shifted parameter
$k + \hv$ is the universal affine denominator or its numerator-dual
shadow.
\end{proposition}

\begin{proof}
\emph{Item (1).} Sugawara's formula for the affine energy--momentum
tensor is derived by requiring that $T(z)$ generate a Virasoro
subalgebra under the affine Kac--Moody OPE~\eqref{eq:affine-km-ope},
with $J^a(z)$ a primary of weight~$1$. The derivation produces the
normalisation factor $1/(2(k + \hv))$, where $2$ comes from the double
pole in the Killing form of the orthonormal basis for $(\ ,\ )$ and
$\hv$ comes from the quadratic Casimir eigenvalue on the adjoint
representation of $\fg$. The composite $k + \hv$ is the SAME
denominator.

\emph{Item (2).} The Sugawara central charge is computed from the
$T(z) T(w)$ OPE. The quartic term in
$:\!J^a J^a\!:\cdot:\!J^b J^b\!:$ contracts pairwise, producing a
double-pole coefficient of $(k \cdot \dim\fg \cdot 2) \cdot
(2(k + \hv))^{-2} \cdot 2 = k\,\dim\fg / (2(k + \hv)^2) \cdot 2\,(k+\hv) = k\,\dim\fg / (k + \hv)$.
(The intermediate factor $2(k + \hv)$ comes from the Killing-form
contraction in the Wick expansion.) Equivalently, one can derive
$c = k\,\dim\fg / (k + \hv)$ from the action of $L_0$ on the vacuum
and the centre of $V_k(\fg)$.

\emph{Item (3).} Direct census entry, landscape\_census.tex:
$\kch^{\mathrm{KM}}(V_k(\fg)) = \dim(\fg)(k + \hv)/(2\hv)$. At $k = 0$
this gives $\dim(\fg)/2 \neq 0$, which is the correct abelian limit.
At $k = -\hv$ this gives $0$, the critical-level
value.

\emph{Item (4).} The KZ connection is derived by requiring horizontal
sections to satisfy the KZ equation, which in turn is derived from
the Ward identity for $V_k(\fg)$-conformal blocks. The coefficient
$1/(k + \hv)$ is forced by the Sugawara formula in item~(1).

\emph{Item (5).} At $k = -\hv$, the denominator-controlled
quantities in items~(1), (2), and (4) degenerate: the Sugawara
$T(z)$ is undefined (division by zero), $c$ is undefined, and the
KZ connection has a pole in~$k$. The trace-form shadow invariant in
item~(3) is the numerator-dual reflection of the same shift: it is
regular and equals zero at the critical point. This is the
Feigin--Frenkel critical level.

In each item the numerator is $k$-linear (or $k$-independent in
items~(1), (4)); the only denominator of $k$ that appears is
$k + \hv$. This is the universality statement.
\end{proof}

\section{Dual Coxeter coincidence via the strange formula}
\label{sec:strange-formula}

We illustrate the universality with the explicit dependence on $\hv$
through the Freudenthal--de~Vries strange formula.

\begin{theorem}[Dual Coxeter coincidence]
\label{thm:h-dual-coxeter-coincidence}
\ClaimStatusProvedHere
Let $\fg$ be simple with rank $r$, dimension $d = \dim\fg$, dual
Coxeter number $\hv$, and half-sum of positive roots $\rho$. The
Freudenthal--de~Vries strange formula
\begin{equation}\label{eq:strange}
 \frac{(\rho, \rho)}{2\hv} \;=\; \frac{d}{24}
\end{equation}
(in the normalised invariant form) implies the identity
\begin{equation}\label{eq:kappa-vs-c}
 \kch(V_k(\fg))
 \;=\; \frac{d(k + \hv)}{2\hv}
 \;=\; \frac{12(\rho, \rho)(k + \hv)}{d \cdot 2\hv}
 \cdot \frac{d^2}{12(\rho,\rho)}
 \;=\; \frac{c(V_k(\fg))}{2}
 \cdot \frac{(k + \hv)^2}{k\,\hv}
\end{equation}
at $k \neq 0, -\hv$. In particular the ratio
$\kch(V_k(\fg))/c(V_k(\fg)) = (k + \hv)^2 / (2k\,\hv)$ is a rational
function of $k$ with its only pole at $k = 0$ (abelian limit) and
its only zero at $k = -\hv$ (critical limit). Explicit evaluations:
\begin{center}
\begin{tabular}{lccccc}
\toprule
$\fg$ & $\fsl_2$ & $\fsl_3$ & $\mathfrak{so}_5$ & $G_2$ & $E_8$ \\
\midrule
$\hv$ & $2$ & $3$ & $3$ & $4$ & $30$ \\
$d$ & $3$ & $8$ & $10$ & $14$ & $248$ \\
$\kch$ at $k = 1$ & $9/4$ & $16/3$ & $20/3$ & $35/4$ & $248 \cdot 31 / 60$ \\
$c$ at $k = 1$ & $1$ & $2$ & $5/2$ & $14/5$ & $248 / 31$ \\
\bottomrule
\end{tabular}
\end{center}
\end{theorem}

\begin{proof}
Equation~\eqref{eq:strange} is the Freudenthal--de~Vries strange
formula; see Kac, \emph{Infinite-Dimensional Lie Algebras},
Ch.~12. Substituting into
$\kch(V_k(\fg)) = d(k + \hv)/(2\hv)$ and
$c(V_k(\fg)) = kd/(k + \hv)$ produces the identity~\eqref{eq:kappa-vs-c}
after algebra.

For the explicit evaluations, at $k = 1$:
\begin{itemize}
 \item $\fsl_2$: $\hv = 2$, $d = 3$, $\kch = 3 \cdot 3 / 4 = 9/4$,
 $c = 3/3 = 1$.
 \item $\fsl_3$: $\hv = 3$, $d = 8$, $\kch = 8 \cdot 4 / 6 = 32/6$,
 $c = 8/4 = 2$.
 \item $\mathfrak{so}_5$: $\hv = 3$, $d = 10$,
 $\kch = 10 \cdot 4 / 6 = 20/3$, $c = 10/4 = 5/2$.
 \item $G_2$: $\hv = 4$, $d = 14$,
 $\kch = 14 \cdot 5 / 8 = 35/4$, $c = 14/5$.
 \item $E_8$: $\hv = 30$, $d = 248$,
 $\kch = 248 \cdot 31 / 60$, $c = 248/31$.
\end{itemize}
Each case exhibits the denominator $k + \hv$ (with $k = 1$) and the
compensating numerator factor $d = \dim\fg$.
\end{proof}

\begin{remark}[Reading of the strange-formula identity]
\label{rem:strange-formula-reading}
The identity~\eqref{eq:kappa-vs-c} is not a coincidence; it is the
shadow of the Sugawara construction itself. The ratio
$\kch/c = (k + \hv)^2/(2k\hv)$ records the two independent places
where the Sugawara denominator enters the chiral datum: once in the
normalisation of $T(z)$ (the $c$ side) and once in the level
parametrisation of the shadow invariant (the $\kch$ side). The strange
formula aligns the $\hv$-dependence of both.
\end{remark}

\section{The trace-form / KZ convention bridge}
\label{sec:bridge}

The bridge between the two $r$-matrix presentations takes the
following sharp form.

\begin{theorem}[Trace--KZ convention bridge]
\label{thm:trace-KZ-convention-bridge}
\ClaimStatusProvedHere
Let $\fg$ be simple and $k \neq -\hv$. The two classical $r$-matrix
presentations $\rtr(z) = k\,\Omegatr / z$ and
$\rKZ(z) = \OmegaKZ / ((k + \hv) z)$ are NOT equal as
$\fg \otimes \fg$-valued rational functions of $k$ and $z$: at $k = 0$
the former vanishes while the latter equals $\OmegaKZ / (\hv\,z) \neq 0$
for non-abelian $\fg$. With the normalisation
$\OmegaKZ=\Omegatr$, their ratio is
\begin{equation}\label{eq:trace-kz-ratio}
 \frac{\rtr(z)}{\rKZ(z)} = k(k+\hv).
\end{equation}
Consequently the two tensor formulas coincide only at the roots of
$k(k+\hv)=1$ in this normalisation, or more generally at the roots of
$k(k+\hv)=\lambda$ if one replaces $\OmegaKZ$ by
$\lambda\Omegatr$ using a fixed invariant-form rescaling. No such
rescaling gives equality for generic~$k$.

The bridge is therefore not a tensor identity. It is the following
commuting dictionary from the affine OPE:
\[
 \text{central double pole}
 \longmapsto
 k\,\Omegatr/z,
 \qquad
 \text{Sugawara Ward identity}
 \longmapsto
 \OmegaKZ/((k+\hv)z).
\]
Both arrows start from the same affine chiral datum, but they record
different coefficients of that datum.
\end{theorem}

\begin{proof}
The inequality of $\rtr(z)$ and $\rKZ(z)$ as rational functions follows
from evaluating at $k = 0$: $\rtr(z)|_{k = 0} = 0$, whereas
$\rKZ(z)|_{k = 0} = \OmegaKZ/(\hv\,z)$, which is non-zero for
non-abelian $\fg$ (the Lie bracket persists through $\OmegaKZ$).

The normalisation statement is
$\OmegaKZ=\Omegatr=2\hv\,\OmegaKil$, proved in
\eqref{eq:casimir-two-forms}--\eqref{eq:kz-casimir-normalization}.
Substituting it into the trace-form residue gives
\[
 \rtr(z)=\frac{k\,\OmegaKZ}{z}.
\]
Since $\rKZ(z)=\OmegaKZ/((k+\hv)z)$, division gives
\eqref{eq:trace-kz-ratio}. Equality after a fixed rescaling
$\OmegaKZ=\lambda\Omegatr$ would require $k(k+\hv)=\lambda$.
The left side is a non-constant quadratic polynomial in~$k$, so no
constant $\lambda$ gives equality for generic~$k$. For a chosen
$\lambda$, equality occurs exactly at the two roots
\[
 k=\frac{-\hv\pm\sqrt{(\hv)^2+4\lambda}}{2}
\]
when the discriminant is non-zero.

Finally, the source of the two coefficients is different. The
trace-form coefficient $k$ is read directly from the double pole in
the affine OPE. The KZ coefficient $(k+\hv)^{-1}$ is read from the
Sugawara Ward identity, whose derivation is given in the proof of
Theorem~\ref{thm:level-shift-universality}. This proves the dictionary
and excludes the false tensor identity.
\end{proof}

\begin{remark}[Convention-separated universality]
\label{rem:chern-weil-474}
The naive rational identity
$k\,\Omegatr = \OmegaKZ/(k + \hv)$ fails at $k = 0$ and therefore is
NOT an identity in rational functions of $k$. The structural positive
statement is supplied by
Theorems~\ref{thm:level-shift-universality} and
\ref{thm:trace-KZ-convention-bridge}: they identify the denominator
$k + \hv$ as the universal Sugawara parameter while keeping the
trace-form central residue $k\,\Omegatr/z$ separate from the KZ
connection coefficient $\OmegaKZ/((k+\hv)z)$. The same separation is
the holographic datum: central residues and Ward-identity coefficients
share an affine OPE source but live in different normalisations.
\end{remark}

\section{Universal healing}
\label{sec:ap126-healing}

The healing closes at the universal level.

\begin{corollary}[universal healing]
\label{cor:ap126-healing-universal}
\ClaimStatusProvedHere
Every $r$-matrix write for $V_k(\fg)$ in the programme must declare
its convention and exhibit the corresponding $k = 0$ check, as follows.
\begin{enumerate}
 \item \emph{Trace-form write.} The $r$-matrix is written with an
 EXPLICIT $k$ numerator prefix:
 $r(z) = k\,\Omegatr / z$. The $k = 0$ check yields $r = 0$
 (abelian limit). This is the convention adopted in
 Definition~\ref{def:rmatrix-convention}(1).
 \item \emph{KZ write.} The $r$-matrix is written with an EXPLICIT
 $(k + \hv)$ denominator:
 $r(z) = \OmegaKZ / ((k + \hv)\,z)$. The $k = 0$ check yields
 $r = \OmegaKZ / (\hv\,z) \neq 0$ for non-abelian $\fg$, reflecting
 the persistence of the Lie bracket at $k = 0$ in the KZ
 normalisation. This is the convention adopted in
 Definition~\ref{def:rmatrix-convention}(2).
\end{enumerate}
A bare write $r(z) = \Omega/z$ with no level prefix and no
$(k + \hv)$ denominator is a violation; by
Theorem~\ref{thm:level-shift-universality} and
Theorem~\ref{thm:trace-KZ-convention-bridge}, every bare write admits
a UNIQUE augmentation into either form~(1) or~(2), determined by the
convention the author has chosen (or is writing within). The healing
is universal: the level shift $k \mapsto k + \hv$ supplies the KZ
denominator, while the affine OPE supplies the trace-form numerator.
\end{corollary}

\begin{proof}
The trace-form augmentation is the assignment
$r(z) = \Omega/z \mapsto k\,\Omegatr/z$, which is well-defined once
$\Omega$ is identified as $\Omegatr$. The KZ augmentation is
$r(z) = \Omega/z \mapsto \OmegaKZ/((k + \hv) z)$, which is well-defined
once $\Omega$ is identified as $\OmegaKZ$. By
Theorem~\ref{thm:trace-KZ-convention-bridge} the two augmentations are
not tensor-equal for generic~$k$; they are the two conventionally
separate coefficients extracted from the same affine chiral datum. By
Theorem~\ref{thm:level-shift-universality} the denominator in the KZ
augmentation is uniquely forced by the Sugawara primary-field
condition. Hence every bare write admits a unique choice of
augmentation per convention, and the resulting expression is checkable
against its $k=0$ boundary.
\end{proof}

\begin{remark}[downstream]
\label{rem:ap141-downstream}
The level-shift universality also heals (level prefix omission):
once Corollary~\ref{cor:ap126-healing-universal} is applied, every
$r$-matrix write carries either an explicit $k$ numerator (trace
form) or an explicit $(k + \hv)$ denominator (KZ), so the level is
never implicit. The grep check
\[
 \mathtt{r\backslash(z\backslash)\backslash s*=\backslash s*\backslash\backslash Omega\backslash s*/\backslash s*z}
\]
for bare $\Omega/z$ writes is then a complete detector of the
violation class.
\end{remark}

\section{Independent verification}
\label{sec:hz-iv-level-shift}

The four proved-here theorems of this chapter have independent
verification witnesses recorded in
\texttt{compute/tests/test\_chern\_weil\_level\_shift.py}.

\begin{itemize}
 \item \textbf{thm:level-shift-universality}, derived from: this
 chapter's conditions (i)--(iii); verified against: Kac
 (\emph{Infinite-Dim.\ Lie Algebras}, Ch.~12) Sugawara formula
 \`a la Knizhnik--Zamolodchikov; disjoint because Kac's derivation
 of the Sugawara denominator uses the quadratic Casimir eigenvalue
 on the adjoint, not the $r$-matrix CYBE.
 \item \textbf{prop:shift-appears-universally}, derived from:
 landscape\_census.tex line for $\kch^{\mathrm{KM}}$; verified
 against: Di~Francesco--Mathieu--S\'en\'echal, \emph{Conformal
 Field Theory} \S15.3 (KZ equations); disjoint because DMS derive
 the KZ connection from the Ward identity for conformal blocks,
 not from the landscape census.
 \item \textbf{thm:h-dual-coxeter-coincidence}, derived from:
 Freudenthal--de~Vries strange formula; verified against:
 Wendlandt, \emph{The restricted Yangian and Yangian doubles}
 (arXiv:2105.01889), \S3 (trace-form normalisation of the Yangian
 $r$-matrix); disjoint because Wendlandt derives the trace-form
 prefactor from Yangian Hopf algebra axioms, not from the strange
 formula.
 \item \textbf{thm:trace-KZ-convention-bridge}, derived from: this
 chapter's explicit computation of $\rtr / \rKZ$; verified against:
 explicit SymPy verification at $\fg = \fsl_2$, $k = 1$, $c = 1$
 (free-field $\widehat{\fsl}_2{}_1 = $ lattice VOA); disjoint
 because the SymPy test computes the $r$-matrices from the OPE
 on the free-field realisation, not from the change-of-basis
 identity.
\end{itemize}

\section{False coincidence isolated}
\label{sec:fp-cache-entry}

The possible confusion is the following:
\emph{trace-form $r$-matrix $k \Omega/z$ and KZ $r$-matrix
$\Omega/((k + \hv)z)$ are the same tensor after a hidden Casimir
rescaling.} Theorem~\ref{thm:level-shift-universality} proves the
opposite. The denominator $k+\hv$ is the universal Sugawara
denominator; the trace-form coefficient $k$ and the KZ coefficient
$(k+\hv)^{-1}$ are different extractions from the same affine OPE.
The bridge of Theorem~\ref{thm:trace-KZ-convention-bridge} is a
dictionary, not a rational identity.

\section*{Closing: the level shift as a Platonic invariant}

The assignment $k \mapsto k + \hv$ is not a convention-dependent
algebraic manipulation; it is the universal Sugawara parameter of
the affine Kac--Moody chiral algebra. Every quantity attached to
$V_k(\fg)$ that cares about the Sugawara normalisation depends on
$k$ through $k + \hv$ alone. The trace-form and KZ presentations of
$r(z)$ remain convention-separated: one reads the central double pole,
the other reads the Sugawara Ward identity. The critical-level locus
$k = -\hv$ is the zero of this universal denominator, and
Feigin--Frenkel localisation is the chart of the chiral datum adapted
to that locus. The convention discipline is thereby made universal in
Corollary~\ref{cor:ap126-healing-universal}: every bare $r$-matrix
write admits a unique augmentation per convention, and each
augmentation carries its own $k=0$ check.
